\documentclass[book.tex]{subfiles}

\begin{document}
\chapter{Spinors}
\label{chap:hilbert}
\noindent\rule{11cm}{0.4pt} \\
\textbf{Prerequisites:} \autoref{chap:math_basics}, \autoref{chap:vectors}, \autoref{chap:manifold} \\
\textbf{Difficulty Level:} ** \\
\noindent\rule{11cm}{0.4pt}
\vspace{5mm}

Spinors are foundational mathematical objects that prove useful for modeling certain types of particles mathematically (fermions in particular). Mathematically, spinors are defined as fundamental representations of Clifford algebras so we begin this chapter by introducing the formalism of Clifford algebras.

\section{Clifford Algebra}

Let $V$ be a vector space over the complex numbers $\mathbb{C}$. Let $g: V \times V \to \mathbb{C}$ be a quadratic form (roughly, a quadratic form can be thought of a quadratic polynomial defined on $V$). The Clifford algebra $\rm{Cl}(V, g)$ is the algebra generated on $V$ that satisfies the following anticommutation relationship.

\begin{align}
    xy + yx = 2g(x, y)
\end{align}

In the case $V = \mathbb{C}^n$, the bilinear form $g$ can be defined by
\begin{align}
    g(x, y) &= x^T y \\
    &= x_1 y_1 + \dotsc + x_n y_n
\end{align}

Another common example of a Clifford algebra is denoted $\rm{Cl}_{p, q}(\mathbb{R})$. This algebra is built from $p+q$ mutually orthogonal basis vectors using addition and multiplication. These algebras satisfy the following multiplication rule
\begin{align}
    e_i e_j &= \begin{cases}
    +1 & i=j, i \in (1,\dotsc, p) \\
    -1 & i=j, i \in (p+1,\dotsc,p+q) \\
    -e_je_i & i \neq j
    \end{cases}
\end{align}

A representation of a Clifford algebra represents these abstract elements with specific matrices of complex numbers. In the next section, we will learn about a specific such representation.

\section{Gamma Matrices}

The Gamma matrices $\{\gamma^0, \gamma^1, \gamma^2, \gamma^3\}$ (also called the Dirac matrices) are a specific set of $4 \times 4$ matrices that generate a representation of $\rm{Cl}_{1,3}(\mathbb{R})$. The gamma matrices are defined as
\begin{align}
    \gamma^0 &= \begin{pmatrix}
    1 & 0 & 0 & 0 \\
    0 & 1 & 0 & 0 \\
    0 & 0 & -1 & 0 \\
    0 & 0 & 0 & -1 
    \end{pmatrix} \\
    \gamma^1 &= \begin{pmatrix}
    0 & 0 & 0 & 1 \\
    0 & 0 & 1 & 0 \\
    0 & -1 & 0 & 0 \\
    -1 & 0 & 0 & 0 
    \end{pmatrix} \\
    \gamma^2 &= \begin{pmatrix}
    0 & 0 & 0 & -i \\
    0 & 0 & i & 0 \\
    0 & i & 0 & 0 \\
    -i & 0 & 0 & 0 
    \end{pmatrix} \\
    \gamma^3 &= \begin{pmatrix}
    0 & 0 & 1 & 0 \\
    0 & 0 & 0 & -1 \\
    -1 & 0 & 0 & 0 \\
    0 & 1 & 0 & 0 
    \end{pmatrix} \\
\end{align}
These matrices satisfy the following anticommutation relationship
\begin{align}
    \{\gamma^\mu, \gamma^\nu\} &= \gamma^{\mu} \gamma^\nu + \gamma^\nu \gamma^\mu \\
    &= 2 \eta^{\mu\nu} I_4
\end{align}
Here $\eta^{\mu\nu}$ denotes an index in the Minkowski metric
\begin{align}
    \eta &= \begin{pmatrix}
    1 & 0 & 0 & 0 \\
    0 & -1 & 0 & 0 \\
    0 & 0 & -1 & 0 \\
    0 & 0 & 0 & -1 
    \end{pmatrix}
\end{align}
and $I_4$ is the identity matrix
\begin{align}
    I_4 &= \begin{pmatrix}
    1 & 0 & 0 & 0 \\
    0 & 1 & 0 & 0 \\
    0 & 0 & 1 & 0 \\
    0 & 0 & 0 & 1 
    \end{pmatrix}
\end{align}

\section{Weyl Spinor}

Let $M$ be a pseudo-Riemannian manifold with dimension $(p, q)$. Recall that for $x \in M$, $T_x M$ denotes the tangent space at $x$, which is a vector space. We can construct a Clifford algebra on this vector space using the metric tensor $g: T_x M \times T_x M \to \mathbb{R}$ as the quadratic form. Let $\{e_i\}$ form a vector space basis on $T_x M$. Then let $\rm{Cl}(p, q)$ denote the Clifford algebra constructed from this basis. The Weyl spinors are defined as

\begin{align}
    w_j &= \frac{1}{\sqrt{2}}(e_{2j} + i e_{2j + 1}) \\
    w_j^* &= \frac{1}{\sqrt{2}}(e_{2j} - i e_{2j + 1}) 
\end{align}

%\section{Bispinors}

%\textcolor{red}{TODO: Fill out}

\section{Exercises}
\begin{enumerate}
    \item Verify the relationship 
    \begin{align}
        \{\gamma^0, \gamma^0\} &= 2 \eta^{0,0} I_4
    \end{align}
\end{enumerate}
\end{document}
